\hypertarget{chapter-16-custom-document-classes}{%
\chapter{Custom Document Classes}\label{chapter-16-custom-document-classes}}

Every LaTeX document begins with \texttt{\textbackslash{}documentclass}, and for most users the choice is limited to whichever standard class, \texttt{article}, \texttt{report}, \texttt{book}, happens to be closest to what they need, patched afterward with enough preamble packages and redefinitions to approximate the actual desired result. For a single document, this is a reasonable way to work. For a research program producing many related books over years, each one wanting the same title-page design, the same front-matter structure, the same theorem styles and notation conventions, patching a standard class identically in every new document's preamble is both repetitive and fragile: any improvement to the shared conventions has to be copied by hand into every existing document's preamble to take effect there too, and inevitably some copy falls out of sync with the rest.

A custom document class solves this the way a shared library solves the equivalent problem in software: the conventions live in one place, a \texttt{.cls} file, and every document that uses them does so by loading that one file, so an improvement made once is available everywhere the class is used, and every document built with a given version of the class is guaranteed to share its conventions exactly, rather than approximately.

\hypertarget{anatomy-of-a-.cls-file}{%
\section{\texorpdfstring{Anatomy of a \texttt{.cls} File}{Anatomy of a .cls File}}\label{anatomy-of-a-.cls-file}}

A document class file is not fundamentally different from any other body of LaTeX macros; what distinguishes it is that it is loaded through \texttt{\textbackslash{}documentclass} rather than \texttt{\textbackslash{}usepackage}, and by convention it is responsible for loading a base class underneath it, typically \texttt{book} or \texttt{report} for a project of the scale this book has been discussing, via \texttt{\textbackslash{}LoadClass}. A minimal custom class might do nothing more than load \texttt{book}, set the page geometry, load \texttt{fontspec} and set the project's standard typeface, and load the packages (\texttt{amsthm}, \texttt{mathtools}, \texttt{hyperref}) every document in the project needs regardless of content, replacing what would otherwise be a preamble block copied identically into every book's top-level file with a single \texttt{\textbackslash{}documentclass\{myresearchclass\}} line.

Beyond this minimal case, a class file can define new commands specific to the project's document type, redefine how existing structural elements like chapter headings or the table of contents are typeset, and declare class-specific options, \texttt{\textbackslash{}documentclass{[}draft{]}\{myresearchclass\}} versus \texttt{\textbackslash{}documentclass{[}final{]}\{myresearchclass\}}, that change the class's behavior at load time, using \texttt{\textbackslash{}DeclareOption} and \texttt{\textbackslash{}ProcessOptions} to parse whatever options a specific document invocation passes in. This options mechanism is what lets a single class file serve every book in a series while still allowing per-document variation, a draft watermark, a different page size for a print edition, without those variations requiring a different class file for every case that needs them.

\hypertarget{enforcing-front-matter-at-the-class-level}{%
\section{Enforcing Front Matter at the Class Level}\label{enforcing-front-matter-at-the-class-level}}

The most direct payoff of a custom class, for a project maintaining consistent title pages and front matter across many books as discussed later in this part, is the ability to require and structure metadata at the class level rather than leaving it to each document's own ad hoc preamble. A class can define its own commands, \texttt{\textbackslash{}bookauthor}, \texttt{\textbackslash{}bookaffiliation}, \texttt{\textbackslash{}bookdate}, distinct from the generic \texttt{\textbackslash{}author} and \texttt{\textbackslash{}date} the base class already provides, giving the class author full control over exactly what fields a title page requires and exactly how they are formatted, rather than working within whatever \texttt{\textbackslash{}maketitle} the base class happens to already implement.

Where a base class's \texttt{\textbackslash{}maketitle} is difficult to adapt without extensive redefinition, because it was not designed with a particular title-page layout in mind, a custom class can simply define its own title-page command entirely, \texttt{\textbackslash{}makebooktitle}, built from scratch to lay out exactly the fields the project's title pages should contain, in exactly the arrangement the project has settled on, with any field the class considers mandatory checked for presence, and an author-facing warning or error raised at compile time if a required field, an affiliation, a required copyright notice, is missing rather than that omission only being caught by a human proofreading the rendered title page after the fact.

This is worth treating as a genuine specification, not merely a convenience: a class-level requirement that every document declares its title, author, and date before \texttt{\textbackslash{}makebooktitle} can run is a guarantee, enforced by the compiler itself, that no document built with this class can accidentally ship without that information, which is a stronger guarantee than any style guide or personal checklist can provide on its own.

\hypertarget{one-class-many-books}{%
\section{One Class, Many Books}\label{one-class-many-books}}

For a project maintaining several related book-length works, a single shared class file, versioned in its own repository or as a shared submodule included by each book's own repository, is the mechanism that keeps title-page design, notation conventions, and theorem styles genuinely identical across the whole series rather than merely similar. A change to the class, adjusting the title-page layout, adding a new theorem style, updating the standard bibliography formatting, propagates to every book the next time that book is rebuilt against the updated class, without requiring the change to be manually reapplied to each book's own preamble.

This does introduce a dependency-management question a single-document project does not have: which version of the shared class does each book in the series actually depend on, and what happens when the class changes in a way that would alter how an already-published book renders. The same discipline Chapter 4 recommended for reproducibility generally applies here directly: pinning each book's build to a specific, recorded version of the shared class, whether through a Git submodule at a fixed commit or an explicit version note in the book's own build documentation, so that rebuilding an earlier book from source still produces the same output it originally did, even after the shared class has continued to evolve for newer books in the series.

The next chapter turns to what a custom class like this is often built specifically to support: professional title pages, front and back matter conventions, and the print-ready output requirements a book intended for real publication, rather than only PDF distribution, actually needs.
